% Part: incompleteness
% Chapter: incompleteness-provability
% Section: second-incompleteness-thm

\documentclass[../../../include/open-logic-section]{subfiles}

\begin{document}

\olfileid[id]{inc}{inp}{2in}

\olsection{Teorema Ketaklengkapan Kedua}

Bagaimana kita dapat mengungkapkan pernyataan bahwa $\Th{PA}$ tidak
membuktikan konsistensinya sendiri? Menyatakan bahwa $\Th{PA}$ tak konsisten
sama artinya dengan menyatakan bahwa $\Th{PA} \Proves \eq[0][1]$. Jadi,
kita dapat memilih pernyataan konsistensi $\OCon[\Th{PA}]$ sebagai
!!{sentence} $\lnot \OProv[\Th{PA}](\gn{\eq[0][1]})$, dan teorema berikut
memberi hasil yang kita inginkan:

\begin{thm}
\ollabel{thm:second-incompleteness}
Dengan mengasumsikan $\Th{PA}$ konsisten, $\Th{PA}$ tidak !!{derive}
$\OCon[\Th{PA}]$.
\end{thm}

Penting diperhatikan bahwa teorema ini bergantung pada representasi khusus
$\OCon[\Th{PA}]$ (yakni representasi khusus $\OProv[\Th{PA}](y)$). Yang
akan kita gunakan hanyalah bahwa representasi $\OProv[\Th{PA}](y)$ memenuhi
ketiga syarat !!{derivability}, sehingga teorema ini berlaku lebih umum bagi
setiap teori dengan predikat !!{derivability} yang memiliki sifat-sifat
tersebut.

Sketsa argumen G\"odel layak dibaca karena teorema tersebut kemudian muncul
seperti penutup jenaka yang tepat. Argumennya berjalan sebagai berikut.
Misalkan $!G_\Th{PA}$ adalah kalimat G\"odel yang kita bangun dalam
pembuktian \olref[1in]{thm:first-incompleteness}. Kita telah menunjukkan,
``Jika $\Th{PA}$ konsisten, maka $\Th{PA}$ tidak !!{derive}
$!G_\Th{PA}$.'' Jika kita memformalkan hal ini \emph{di dalam} $\Th{PA}$,
kita memperoleh pembuktian atas
\[
\OCon[\Th{PA}] \lif \lnot \OProv[\Th{PA}](\gn{!G_\Th{PA}}).
\]
Sekarang andaikan $\Th{PA}$ !!{derive}s $\OCon[\Th{PA}]$. Maka teori itu
!!{derive}s $\lnot \OProv[\Th{PA}](\gn{!G_\Th{PA}})$. Namun, karena
$!G_\Th{PA}$ adalah kalimat G\"odel, formula ini ekuivalen dengan
$!G_\Th{PA}$. Jadi, $\Th{PA}$ !!{derive}s $!G_\Th{PA}$.

Namun, kita tahu bahwa jika $\Th{PA}$ konsisten, teori itu tidak !!{derive}
$!G_\Th{PA}$!{} Jadi, jika $\Th{PA}$ konsisten, teori itu tidak mungkin
!!{derive} $\OCon[\Th{PA}]$.

Untuk membuat argumen tersebut lebih cermat, kita akan memilih
$!G_\Th{PA}$ sebagai kalimat G\"odel bagi~$\Th{PA}$ dan menggunakan
syarat !!{derivability} (P1)--(P3) untuk menunjukkan bahwa $\Th{PA}$
!!{derive}s $\OCon[\Th{PA}] \lif !G_\Th{PA}$. Hal ini akan menunjukkan
bahwa $\Th{PA}$ tidak !!{derive} $\OCon[\Th{PA}]$. Berikut sketsa
pembuktiannya di dalam~$\Th{PA}$. (Demi kesederhanaan, kita menghilangkan
subskrip $\Th{PA}$.)
\begin{align}
& !G \liff \lnot \OProv(\gn{!G}) \ollabel{G2-1}\\
& \qquad\text{$!G$ adalah kalimat G\"odel}\notag \\
& !G \lif \lnot \OProv(\gn{!G}) \ollabel{G2-2}\\
  & \qquad\text{dari \olref{G2-1}} \notag\\
& !G \lif
  (\OProv(\gn{!G}) \lif \lfalse) \ollabel{G2-3}\\
  & \qquad\text{dari \olref{G2-2} berdasarkan logika}\notag\\
& \OProv(\gn{
    !G \lif
    (\OProv(\gn{!G}) \lif \lfalse)
  }) \ollabel{G2-4}\\
  & \qquad\text{dari \olref{G2-3} berdasarkan syarat P1} \notag\\
& \OProv(\gn{!G}) \lif
  \OProv(\gn{
    (\OProv(\gn{!G}) \lif \lfalse)
    }) \ollabel{G2-5}\\
  & \qquad\text{dari \olref{G2-4} berdasarkan syarat P2} \notag\\
& \OProv(\gn{!G}) \lif (\OProv(\gn{\OProv(\gn{!G})}) \lif \OProv(\gn{\lfalse})) \ollabel{G2-6}\\
  & \qquad\text{dari \olref{G2-5} berdasarkan syarat P2 dan logika} \notag\\
& \OProv(\gn{!G}) \lif
  \OProv(\gn{\OProv(\gn{!G})}) \ollabel{G2-7}\\
   & \qquad\text{berdasarkan P3} \notag\\
& \OProv(\gn{!G}) \lif \OProv(\gn{\lfalse}) \ollabel{G2-8}\\
  & \qquad \text{dari \olref{G2-6} dan \olref{G2-7} berdasarkan logika}\notag\\
& \OCon \lif \lnot \OProv(\gn{!G}) \ollabel{G2-9}\\
  & \qquad\text{kontraposisi \olref{G2-8} dan $\OCon \ident \lnot \OProv(\gn{\lfalse})$}\notag \\
& \OCon \lif !G \notag\\
  & \qquad\text{dari \olref{G2-1} dan \olref{G2-9} berdasarkan logika}\notag
\end{align}
Penggunaan logika di atas hanya melibatkan fakta-fakta dasar logika
proposisional; misalnya, \olref{G2-3} menggunakan
$\Proves \lnot!A \liff (!A\lif \lfalse)$ dan \olref{G2-8} menggunakan
$!A \lif (!B \lif !C), !A \lif !B \Proves !A \lif !C$. Penggunaan
syarat~P2 dalam \olref{G2-5} dan \olref{G2-6} bergantung pada contoh-contoh
P2, yakni $\OProv(\gn{!A \lif !B}) \lif
(\OProv(\gn{!A}) \lif \OProv(\gn{!B}))$. Dalam contoh pertama,
$!A \ident !G$ dan $!B \ident \OProv(\gn{!G}) \lif \lfalse$; dalam
contoh kedua, $!A \ident \OProv(\gn{!G})$ dan $!B \ident \lfalse$.

Versi teorema ketaklengkapan kedua yang lebih abstrak adalah sebagai berikut:

\begin{thm}
\ollabel{thm:second-incompleteness-gen} Misalkan $\Th{T}$ sebarang teori
konsisten dan !!{axiomatized} yang memperluas $\Th{Q}$, dan misalkan
$\OProv[\Th{T}](y)$ sebarang formula yang memenuhi syarat !!{derivability}
P1--P3 bagi~$\Th{T}$. Maka $\Th{T}$ tidak !!{derive}~$\OCon[\Th{T}]$.
\end{thm}

\begin{prob}
Tunjukkan bahwa $\Th{PA}$ !!{derive}s
$!G_{\Th{PA}} \lif \OCon[\Th{PA}]$.
\end{prob}

\begin{digress}
Pesan utamanya ialah bahwa tidak ada teori konsisten yang ``wajar'' bagi
matematika yang dapat !!{derive} pernyataan konsistensinya sendiri. Andaikan
$\Th{T}$ adalah teori matematika yang mencakup $\Th{Q}$ dan penalaran
``finitistik'' Hilbert (apa pun tepatnya yang dimaksud). Maka keseluruhan
$\Th{T}$ tidak dapat !!{derive} pernyataan konsistensi $\Th{T}$; terlebih
lagi, fragmen finitistiknya juga tidak dapat !!{derive} pernyataan
konsistensi~$\Th{T}$. Dalam arti ini, tidak mungkin ada pembuktian
konsistensi finitistik bagi ``seluruh matematika.''

Ada sedikit kelonggaran dalam menafsirkan istilah ``finitistik,'' dan dalam
makalah tahun 1931, G\"odel mengakui kemungkinan bahwa sesuatu yang kita
anggap ``finitistik'' mungkin berada di luar jenis matematika yang hendak
diformalkan Hilbert. Namun, G\"odel bersikap murah hati; dewasa ini sulit
melihat bagaimana kita dapat menemukan sesuatu yang layak disebut finitistik
tetapi tidak dapat diformalkan, misalnya, dalam~$\Th{ZFC}$, yakni teori
himpunan Zermelo--Fraenkel dengan aksioma pilihan.
\end{digress}

\end{document}
